\chapter{Conclusión}

El desarrollo de este proyecto permitió ampliar el sistema de gestión de deportistas construido previamente, incorporando algoritmos basados en la estrategia de \textit{Divide y Vencerás} para resolver problemas de ordenamiento, búsqueda y selección. A diferencia del trabajo anterior, centrado principalmente en algoritmos iterativos clásicos, esta segunda etapa permitió analizar técnicas recursivas más eficientes y observar cómo su comportamiento depende tanto de la complejidad teórica como de las decisiones concretas de implementación. \\

En relación con los algoritmos de ordenamiento, la implementación de \textit{Merge Sort} permitió contar con un método estable en cuanto a complejidad temporal, ya que mantiene un comportamiento $O(n \log n)$ en mejor caso, caso promedio y peor caso. Esto se debe a que el algoritmo divide el arreglo de manera relativamente uniforme y luego combina los subarreglos mediante un proceso de mezcla lineal. En la práctica, esta propiedad se reflejó en un rendimiento consistente frente a distintos escenarios de entrada, lo que confirma su utilidad cuando se requiere un comportamiento predecible. La variante \textit{Merge-Insertion Sort} permitió estudiar el efecto de combinar dos estrategias de ordenamiento. El uso de \textit{Insertion Sort} en subarreglos pequeños busca reducir el costo asociado a llamadas recursivas demasiado finas, aprovechando que dicho algoritmo suele comportarse bien en entradas de tamaño reducido. Por ello, esta variante resulta relevante desde un punto de vista experimental: aunque mantiene la misma complejidad asintótica que \textit{Merge Sort}, puede presentar mejoras prácticas dependiendo del umbral utilizado. Esto demuestra que dos algoritmos con igual complejidad teórica pueden diferir en rendimiento real debido a constantes ocultas, uso de memoria y sobrecarga de implementación. \\

En el caso de \textit{Quick Sort}, los resultados permiten evidenciar con claridad la importancia de la selección del pivote. Las variantes que utilizan el primer o el último elemento como pivote pueden degradarse significativamente cuando la entrada ya se encuentra ordenada o invertida, debido a que generan particiones altamente desbalanceadas. Este comportamiento coincide con el peor caso teórico de \textit{Quick Sort}, cuya complejidad puede alcanzar $O(n^2)$. En cambio, el pivote aleatorio y la mediana de tres buscan reducir la probabilidad de particiones desfavorables, acercando el comportamiento del algoritmo a su caso promedio $O(n \log n)$. De esta forma, el proyecto muestra que \textit{Quick Sort} no depende únicamente de su estructura general, sino también de una decisión aparentemente simple como la elección del pivote. \\

En relación con los algoritmos de búsqueda, los resultados experimentales revelaron diferencias significativas entre las distintas estrategias implementadas. La búsqueda binaria recursiva presentó tiempos consistentemente bajos, confirming su complejidad teórica $O(\log n)$, con un crecimiento prácticamente imperceptible al aumentar el tamaño de entrada. La búsqueda exponencial, aunque mantiene la misma complejidad asintótica, mostró un comportamiento ligeramente superior en términos de tiempos observados, lo que sugiere que las constantes ocultas asociadas a su fase inicial de expansión impactan el rendimiento en la práctica. Por su parte, la búsqueda por interpolación demostró una dependencia crítica respecto a la distribución de los datos: en entradas uniformemente distribuidas alcanzó rendimientos comparables a la búsqueda binaria, pero degradó notablemente cuando los datos presentaban sesgos o agrupamientos, confirmando la importancia de considerar la naturaleza específica del conjunto de datos. \\

La búsqueda binaria por rangos contribuyó con una funcionalidad práctica al sistema, permitiendo obtener subconjuntos de deportistas dentro de un intervalo de puntajes determinado. Los benchmarks mostraron que esta variante mantiene un comportamiento eficiente incluso al expandir el rango de búsqueda, debido a que realiza dos búsquedas binarias para localizar los límites. La implementación resultó útil para ampliar las capacidades del sistema más allá de búsquedas exactas, facilitando consultas de rango que son comunes en aplicaciones de gestión de datos. \\

La implementación de \textit{Quick Select} permitió resolver el problema de encontrar el $k$-ésimo mejor deportista sin ordenar completamente el arreglo. Esta diferencia es importante, ya que ordenar todos los datos puede ser innecesario cuando solo se necesita conocer una posición específica. Al reutilizar la lógica de partición de \textit{Quick Sort}, \textit{Quick Select} permite descartar una parte del arreglo en cada paso recursivo y concentrarse únicamente en el segmento donde puede encontrarse el elemento buscado. En promedio, esto permite alcanzar complejidad $O(n)$, aunque el algoritmo también puede degradarse a $O(n^2)$ si las particiones son muy desbalanceadas. \\

Desde el punto de vista experimental, los \textit{benchmarks} realizados permitieron contrastar la teoría con el comportamiento real de los algoritmos. Las gráficas generadas a partir de los archivos CSV facilitaron observar tendencias de crecimiento, diferencias entre variantes y efectos de los casos de prueba. En particular, los resultados de ordenamiento mostraron que \textit{Merge Sort} y \textit{Merge-Insertion Sort} poseen un rendimiento más estable frente a distintas entradas, mientras que \textit{Quick Sort} puede variar considerablemente según la variante de pivote utilizada. En las búsquedas, los algoritmos basados en división del espacio de búsqueda mantuvieron tiempos reducidos frente al crecimiento de $n$, mientras que la búsqueda por interpolación dependió con mayor fuerza de la distribución de los datos. \\

En conclusión, el proyecto permitió comprobar que la estrategia de \textit{Divide y Vencerás} ofrece herramientas eficientes para trabajar con conjuntos de datos crecientes, especialmente cuando se compara con enfoques iterativos más simples. Sin embargo, también se evidenció que la complejidad asintótica no explica por sí sola todo el comportamiento práctico de un algoritmo. Factores como la elección del pivote, el uso de memoria auxiliar, el umbral de cambio entre algoritmos, la distribución de los datos y el diseño de los casos experimentales influyen directamente en los tiempos observados. Por ello, la selección de un algoritmo adecuado debe considerar tanto su fundamento teórico como el contexto concreto en que será aplicado. \\